\documentclass[11pt,a4paper]{article}
\usepackage[margin=2.2cm]{geometry}
\usepackage{amsmath,amssymb,booktabs,tabularx,microtype}
\usepackage[colorlinks=true,linkcolor=black,urlcolor=blue]{hyperref}
\usepackage{titlesec}
\titleformat{\section}{\large\bfseries\sffamily}{\thesection.}{0.6em}{}
\titleformat{\subsection}{\normalsize\bfseries\sffamily}{}{0em}{}
\setlength{\parskip}{0.4em}\setlength{\parindent}{0pt}
\newcolumntype{L}{>{\raggedright\arraybackslash}X}
\newcommand{\pro}{\textbf{\textsf{+}}~}
\newcommand{\con}{\textbf{\textsf{--}}~}

\title{\vspace{-1.5cm}\sffamily Adaptive masking: six options\vspace{-0.3em}}
\author{\small Working note, 1 September 2026 --- not part of the thesis}
\date{}

\begin{document}
\maketitle
\vspace{-1.2em}

\section*{The problem this addresses}

Measured 2026-09-01: the extractor applies the \emph{same} transform regardless
of how hard the trial is. Across easy-to-hard buckets $\Delta\text{SIR}$ moves only
$+3.80 \to +4.32$~dB and $\Delta\text{SAR}$ only $-17.45 \to -21.41$~dB, while the
word-error outcome swings from $-4.2$ to $+23.1$ points. It invents roughly
9\,\% of its output energy (absolute SAR $+10.34$~dB) and does so everywhere,
which is a good trade on hard trials and a bad one on easy trials.

\subsection*{What the screening sweep found (1 September 2026)}

The inference-time blend $\hat{S}_\alpha = \alpha\hat{S} + (1-\alpha)X$ was swept
over $\alpha \in \{0,0.25,0.5,0.75,1\}$ on 103 two-speaker trials. SDR, SIR and
SAR all moved perfectly monotonically, confirming the blend behaves as predicted.

\textbf{Globally $\alpha=1$ wins} --- there is no interior optimum, so a global
shift toward gentler masking is the wrong direction. \textbf{But the optimum per
difficulty bucket spans the entire range:}

\begin{center}\small
\begin{tabular}{@{}lccccc@{}}
\toprule
difficulty & $\alpha{=}0$ & $0.25$ & $0.5$ & $0.75$ & $\alpha{=}1$ \\
\midrule
easy ($<$25\,\%)      & \textbf{9.6} & 11.2 & 10.3 & 11.7 & 14.2 \\
medium                & 40.5 & \textbf{36.8} & 58.6 & 46.6 & 40.9 \\
hard                  & 81.9 & 80.6 & 82.7 & 79.1 & \textbf{73.4} \\
very hard ($>$100\,\%) & 137.4 & 133.8 & 131.4 & 140.1 & \textbf{113.6} \\
\bottomrule
\end{tabular}
\end{center}

Easy trials want \emph{no filtering}; hard trials want \emph{full filtering}. A
single constant cannot be right, so \textbf{adaptation is what the data
supports.} An oracle picking $\alpha$ per bucket reaches \textbf{56.9\,\%}
against the model's 59.1\,\% --- about \textbf{2.2 points} available, with
52.5\,\% as an unreachable per-trial ceiling.

\emph{Caveat:} the word-error curve is non-monotonic ($\alpha{=}0.5$ is worse
than doing nothing) while the signal measures are not. The blend is provably
linear, so this is $n{=}103$ noise plus transcriber nonlinearity. Trust the
endpoints and the per-bucket pattern, not individual interior values.

\subsection*{Notation}

$X[t,b]$ is the mixture in the time--frequency domain, frame $t$, band $b$;
$m[t,b]$ the complex mask the estimation module predicts; $h[t,b]$ the model's
internal features; $e$ the enrolment. The current model is
\[
\hat{S}[t,b] \;=\; m[t,b]\,X[t,b].
\]
$T$ frames, $B=32$ bands, $7.19$\,M parameters, trained on $4{,}976$ trials.
\textbf{The model is data-limited, not capacity-limited} --- train separation
improved monotonically for all 18 epochs while held-out peaked at epoch~7. Any
option that adds capacity makes the measured problem worse.

\section{Global mix-back constant (no learning)}
\[
\hat{S}_\alpha \;=\; \alpha\,\hat{S} \;+\; (1-\alpha)\,X,
\qquad \alpha \text{ fixed, chosen offline}
\]
\pro Zero training, zero new parameters, zero added latency: one multiply--add
per sample, and output sample $n$ needs only input sample $n$.\\
\pro Every $\alpha$ comes from \emph{one} forward pass, so a whole family of
systems costs nothing but transcription.\\
\con Not a fix. Trades away the benefit on hard trials to stop the harm on easy
ones, using a constant that cannot be right everywhere.\\
\con No gradient --- $\alpha$ is not a model parameter, so the model learns
nothing.

\section{Per-frame learned gate}
\[
\alpha[t] \;=\; \sigma\!\big(w^{\!\top} h[t] + b\big),
\qquad
\hat{S}[t] \;=\; \alpha[t]\,\big(m[t]\,X[t]\big) + \big(1-\alpha[t]\big)X[t]
\]
Gradient is immediate and needs no trick:
\[
\frac{\partial \hat{S}[t]}{\partial \alpha[t]} \;=\; m[t]\,X[t] - X[t],
\qquad
\frac{\partial \alpha[t]}{\partial w} \;=\; \alpha[t]\big(1-\alpha[t]\big)h[t].
\]
\pro Roughly $1$--$2$\,k parameters against $7.19$\,M: capacity risk negligible.\\
\pro Learns \emph{when} to be gentle, which the global constant cannot.\\
\pro Backbone untouched, so the 2026-08-28 architecture freeze survives.\\
\con One decision for the whole spectrum at each frame.\\
\con Still needs a retrain (${\sim}6.2$\,h per arm).

\section{Per-frame, per-band learned gate}
\[
\alpha[t,b] \;=\; \sigma\!\big(w_b^{\!\top} h[t,b] + b_b\big),
\qquad
\hat{S}[t,b] \;=\; \alpha[t,b]\,\big(m[t,b]\,X[t,b]\big) + \big(1-\alpha[t,b]\big)X[t,b]
\]
\pro The band-split architecture already carries per-band features, so this is
architecturally natural rather than bolted on.\\
\pro Can be gentle in the high bands, where masking artefacts live, and
aggressive in the low bands, where speech energy is. That is a distinction no
scalar gate can express.\\
\pro Still only ${\sim}B\,(|h_b|+1)$ parameters --- tens of thousands, not millions.\\
\con Slightly more to defend in the write-up than a scalar.\\
\con Same retrain cost.

\section{Artefact weighting in the loss ($\beta$) --- no architecture change}

The three components of the estimate are mutually orthogonal, so
$L_\textup{pres}$'s denominator is \emph{already} the sum of the two piles:
\[
\big\lVert \hat{s} - \alpha s \big\rVert^2
  \;=\; \big\lVert e_\textup{interf} \big\rVert^2
      + \big\lVert e_\textup{artif} \big\rVert^2
\]
(verified numerically to six decimal places). They therefore cost the same per
unit of energy. One ratio changes that:
\[
L_\textup{sep} \;=\; -10\log_{10}
  \frac{\lVert \alpha s\rVert^2}
       {\lVert e_\textup{interf}\rVert^2
        + \beta\,\lVert e_\textup{artif}\rVert^2
        + \tau\lVert \alpha s\rVert^2}
\]
\pro $\beta=1$ recovers $L_\textup{pres}$ \emph{exactly}, so the control arm is
the existing baseline and every prior run stays comparable.\\
\pro No architecture change, no new parameters, no latency change.\\
\pro Lets the model choose \emph{where} in time and frequency to be gentle, which
a global $\alpha$ cannot.\\
\pro Unlike a separate artefact loss, pass-through can never win it: the
interference term dominates.\\
\con Shifts the global operating point; does \emph{not} make the model adaptive.\\
\con One new hyperparameter, and a 4-arm ablation is ${\sim}25$\,h of GPU time.

\section{Soft-routing mixture of experts}
\[
g \;=\; \operatorname{softmax}\!\big(f(X,e)\big) \in \mathbb{R}^K,
\qquad
\hat{S} \;=\; \sum_{k=1}^{K} g_k \,\big(m_k \, X\big)
\]
Differentiable without any trick, since softmax is smooth and the sum is linear:
\[
\frac{\partial L}{\partial g_k}
  \;=\; \Big\langle \frac{\partial L}{\partial \hat{S}},\; m_k X \Big\rangle .
\]
\pro No estimator bias, no sampling, no Gumbel trick.\\
\pro Genuinely expresses ``different treatment for different cases''.\\
\con Replicating the whole estimator is fatal: $K{=}5$ would take the model from
$7.19$\,M to $15.9$\,M, $+122\,\%$.\\
\pro \textbf{Sharing the trunk changes the picture.} Measured: the estimator is
$2.19$\,M of which \texttt{trunks} is $1.59$\,M ($73\,\%$), \texttt{mask\_heads}
$396$\,k and \texttt{res\_heads} $198$\,k. Replicating only the mask heads costs
$+0.79$\,M at $K{=}3$, i.e.\ \textbf{$+11\,\%$} --- and the separator ($68\,\%$)
and trunks stay shared, so the experts are $K$ read-outs of one representation
rather than $K$ models.\\
\con All $K$ masks are computed every frame, so runtime scales with $K$.\\
\con Averaging masks can yield a mask worse than any individual one.

\section{Variational mixture of experts (the version asked about)}

Treat difficulty as a discrete latent $z \in \{1,\dots,K\}$ with prior $p(z)$
and posterior $q(z\mid X,e)$, and optimise the evidence lower bound
\[
\mathcal{L} \;=\;
  \mathbb{E}_{q(z\mid X,e)}\!\big[\,L_\textup{recon}(\hat{S}_z, S)\,\big]
  \;-\; \lambda\,\mathrm{KL}\big(q(z\mid X,e)\,\Vert\,p(z)\big).
\]
\textbf{For small $K$ the expectation is computed exactly by enumeration}, so no
sampling, no reparameterisation trick and no REINFORCE variance:
\[
\mathcal{L} \;=\; \sum_{k=1}^{K} q_k\,L_\textup{recon}(\hat{S}_k, S)
  \;-\; \lambda \sum_{k=1}^{K} q_k \log \frac{q_k}{p_k}.
\]
\pro \textbf{Mixes the losses, not the masks.} Each expert must be individually
good on the cases assigned to it, avoiding the averaged-mask failure of
soft routing.\\
\pro The KL term prevents the gate collapsing onto a single expert, and
$\lambda$ controls how decisively it specialises.\\
\pro Principled, citable, and a real architecture contribution.\\
\con $K$ full loss evaluations per step as well as $K$ masks: the most expensive
option by a wide margin.\\
\con $\lambda$ and $p(z)$ to choose on top of $K$.\\
\pro The enumeration point is publishable in its own right: most treatments of
discrete latents reach for Gumbel-Softmax or REINFORCE, and for $K$ of 5--10
neither is needed.\\
\con A new architecture is a different thesis from a new metric.

\section*{Summary}

\begin{tabularx}{\textwidth}{@{}lLcccc@{}}
\toprule
Option & Added parameters & Latency & Capacity risk & Build cost & Fixes non-adaptation \\
\midrule
1. Global $\alpha$        & none                & none & none   & hours (no training) & no \\
2. Per-frame gate         & ${\sim}1$--$2$\,k     & none & tiny   & 1 retrain           & partly \\
\textbf{3. Per-band gate} & ${\sim}10$--$50$\,k   & none & small  & 1 retrain           & \textbf{yes} \\
\textbf{4. $\beta$ in loss} & none              & none & none   & 4 retrains (25\,h)  & no \\
5. Soft MoE               & $+11\,\%$ at $K{=}3$, shared trunk & $K\times$ head & moderate & 1 retrain + tuning & yes \\
6. Variational MoE        & $+11\,\%$ at $K{=}3$, shared trunk & $K\times$ head & moderate & largest & yes \\
\bottomrule
\end{tabularx}

\subsection*{Recommendation, after the sweep}

\textbf{The sweep has been run, and it re-orders the list.}

\textbf{Option 4 ($\beta$) is demoted.} Globally $\alpha=1$ already wins, so a
global shift toward gentler masking is evidence-against rather than
evidence-for. It stays a secondary arm --- the sweep is an imperfect proxy, since
a retrained model learns a \emph{different} mask rather than a blended one --- but
it is no longer first.

\textbf{Option 3 (per-band gate) is now the primary M5 experiment.} The optimum
varies from $\alpha=0$ to $\alpha=1$ across difficulty, so a single constant
provably cannot be right. It costs tens of thousands of parameters against
$7.19$\,M, leaves the backbone frozen, adds no latency, and there are $2.2$
points of LCF-WER on the table.

\textbf{And note the identity that links it to options 5 and 6:}
\[
\alpha\,(m \odot X) + (1-\alpha)\,X
   \;=\; \big[\,\alpha m + (1-\alpha)\mathbf{1}\,\big] \odot X
\]
\textbf{the per-band gate IS a two-expert mixture whose second expert is the
identity mask.} So it is not an alternative to the mixture-of-experts idea --- it
is that idea's base case at near-zero cost, and the right way to test the routing
hypothesis before paying for $K>2$.

\textbf{Options 5 and 6 are deferred, not rejected, and the capacity objection
has been withdrawn.} With a shared trunk, $K{=}3$ costs $+11\,\%$, which is
defensible. What defers them is scope: six weeks to freeze and no live-model
measurement taken. The sensible sequence is \emph{build the $K{=}2$ gate, see how
much of the $2.2$ points it captures, and scale $K$ only if it does.} If the
gate saturates at $\alpha=1$ everywhere, that is a cheap negative result that
saves building the larger thing.

\end{document}
